% Options for packages loaded elsewhere
\PassOptionsToPackage{unicode}{hyperref}
\PassOptionsToPackage{hyphens}{url}
\documentclass[
  11pt,
  a4paper,
]{scrartcl}
\usepackage{xcolor}
\usepackage[top=23mm,bottom=23mm,left=24mm,right=24mm]{geometry}
\usepackage{amsmath,amssymb}
\setcounter{secnumdepth}{-\maxdimen} % remove section numbering
\usepackage{iftex}
\ifPDFTeX
  \usepackage[T1]{fontenc}
  \usepackage[utf8]{inputenc}
  \usepackage{textcomp} % provide euro and other symbols
\else % if luatex or xetex
  \usepackage{unicode-math} % this also loads fontspec
  \defaultfontfeatures{Scale=MatchLowercase}
  \defaultfontfeatures[\rmfamily]{Ligatures=TeX,Scale=1}
\fi
\usepackage{lmodern}
\ifPDFTeX\else
  % xetex/luatex font selection
\fi
% Use upquote if available, for straight quotes in verbatim environments
\IfFileExists{upquote.sty}{\usepackage{upquote}}{}
\IfFileExists{microtype.sty}{% use microtype if available
  \usepackage[]{microtype}
  \UseMicrotypeSet[protrusion]{basicmath} % disable protrusion for tt fonts
}{}
\usepackage{setspace}
\makeatletter
\@ifundefined{KOMAClassName}{% if non-KOMA class
  \IfFileExists{parskip.sty}{%
    \usepackage{parskip}
  }{% else
    \setlength{\parindent}{0pt}
    \setlength{\parskip}{6pt plus 2pt minus 1pt}}
}{% if KOMA class
  \KOMAoptions{parskip=half}}
\makeatother
\usepackage{longtable,booktabs,array}
\newcounter{none} % for unnumbered tables
\usepackage{calc} % for calculating minipage widths
% Correct order of tables after \paragraph or \subparagraph
\usepackage{etoolbox}
\makeatletter
\patchcmd\longtable{\par}{\if@noskipsec\mbox{}\fi\par}{}{}
\makeatother
% Allow footnotes in longtable head/foot
\IfFileExists{footnotehyper.sty}{\usepackage{footnotehyper}}{\usepackage{footnote}}
\makesavenoteenv{longtable}
\setlength{\emergencystretch}{3em} % prevent overfull lines
\providecommand{\tightlist}{%
  \setlength{\itemsep}{0pt}\setlength{\parskip}{0pt}}

\usepackage{microtype}
\usepackage{xcolor}
\usepackage{needspace}
\definecolor{linkink}{HTML}{553653}
\AtBeginDocument{\hypersetup{colorlinks=true,linkcolor=linkink,urlcolor=linkink,citecolor=linkink}}
\setlength{\emergencystretch}{3em}
\setlength{\parindent}{0pt}
\setlength{\parskip}{0.42em}
\widowpenalty=10000
\clubpenalty=10000
\displaywidowpenalty=10000
\allowdisplaybreaks[1]
\usepackage{bookmark}
\IfFileExists{xurl.sty}{\usepackage{xurl}}{} % add URL line breaks if available
\urlstyle{same}
\hypersetup{
  pdftitle={Small service requirements determine information freshness near saturation},
  pdfauthor={Anonymous},
  hidelinks,
  pdfcreator={LaTeX via pandoc}}

\title{Small service requirements determine information freshness near
saturation}
\author{Anonymous}
\date{20 October 2026}

\begin{document}
\maketitle

\setstretch{1.04}
\section{Abstract}\label{abstract}

Consider a stationary lossless M/G/1 queue under egalitarian processor
sharing. Every arrival carries its generation timestamp, and the monitor
retains the greatest timestamp delivered so far. If the service CDF
satisfies \(G(x)\sim cx^\beta\) at zero and has finite mean \(m\), we
prove that the stationary information age, multiplied by
\((1-\rho)^{\beta/(\beta+1)}\), converges to a product of powers of two
independent exponential variables. Every positive real moment converges.
More generally, the same law holds for every service CDF regularly
varying at zero with positive index, with an explicit integral-defined
normalization that also captures slowly varying factors. The proof uses
an exact survival product in a queue with immortal fresh arrivals, a
monotone coupling to bound old departures using only stationarity, and a
uniform exponential tail estimate. No finite service second moment is
required. Queues with identical stationary population distributions can
therefore have different information-age scales. For exponential
service, the exact mean-age coefficient is \(\pi/(2\sqrt2)\) in units of
\(1/(\mu\sqrt{1-\rho})\). Stochastic bounds establish the valid
components of the inspected arXiv version of a 2023 conjecture, and the
limit replaces both printed near-saturation bounds. Neither
processor-sharing insensitivity nor the product distribution is claimed
as new.

\textbf{Keywords:} age of information; processor sharing; heavy traffic;
regular variation; stationary queues; moment convergence.

\section{1. Problem and relation to previous
work}\label{problem-and-relation-to-previous-work}

Age of information measures how old the freshest delivered observation
is. Out-of-order delivery matters: a packet may still occupy the server
even after another packet with a later timestamp has completed. Equal
processor sharing offers fresh packets immediate access to service,
although most of the capacity of a nearly saturated queue is shared with
older work. The question is how quickly stationary age grows as the
arrival rate approaches the stability boundary.

Liu, Huang, Li and Ji (2021) studied processor sharing among competing
age-based scheduling policies and reported simulations with exponential
and Weibull service requirements. Gandarias, Doncel and Assaad (2023)
analysed finite-capacity processor-sharing systems and stated a
conjecture for the mean age of the infinite-buffer M/M/1 system.
Conjecture 1, equations (19)--(20), in their inspected arXiv manuscript
suggests a square-root upper growth scale and leaves the mean correction
undetermined. Its two printed near-saturation bounds are also
algebraically incompatible; noticing that inconsistency alone does not
determine the correct asymptotics.

This paper proves a broader result. For a fixed service distribution
with \(G(x)\sim cx^\beta\) near zero, the age exponent is
\(\beta/(\beta+1)\). The entire limiting law and all positive moments
are identified. The special case \(\beta=1\) supplies the exact
exponential-service constant and, together with separate stochastic
comparisons, settles every mathematically consistent assertion in the
inspected conjecture.

The stationary M/G/1 processor-sharing population is geometric and
insensitive to the service distribution beyond its mean. This classical
fact is an input to the proof, obtained here from Zachary (2007), with
wider network background in Bonald and Proutière (2002). It does not
imply insensitivity of freshness. The completion of one sufficiently
short new job refreshes the monitor, so the lower end of the service
distribution governs the age scale. Section 3 makes this explanation
rigorous without a fluid-limit assumption or a service second-moment
condition.

\textbf{Guide to the argument.} Theorem 1 gives an explicit formula for
the familiar power-law case. Its proof has three distinct tasks: express
age as a first-completion event, show that the large initial population
changes negligibly on the shorter freshness time scale, and control the
tail strongly enough to pass to moments. Section 3.6 then proves the
regularly varying extension, which permits logarithmic corrections to a
power law. Readers primarily interested in consequences can proceed from
the theorems to the equal-mean examples in Section 4 and the
exponential-service comparison in Section 5. The numerical checks in
Section 6 test representations and illustrate finite-load bias; they are
not premises of the proofs.

\section{2. Model and main theorem}\label{model-and-main-theorem}

Packets arrive according to a Poisson process of rate
\(\lambda_\rho=\rho/m\), where \(0<\rho<1\). Their independent and
identically distributed service requirements have CDF \(G\),
independently of arrivals. Total service capacity is one unit of work
per unit time: with \(n\) packets present, each receives work at rate
\(1/n\). The buffer is infinite. Every packet remains until completion,
including packets that have become obsolete.

Write \(a_j\) and \(d_j\) for generation/arrival and departure times.
The stationary age process is

\[
A_\rho(t)=t-\sup\{a_j:d_j\leq t\}.
\]

Thus an old completion does not overwrite a fresher timestamp. This is
the convention implemented by the fake-update transition in Appendix A,
Table I, of Gandarias et al.~(2023). Write \(A_\rho\) for the age at a
stationary observation time. We assume the queue is in its standard
stationary regime and vary only the Poisson rate as \(\rho\uparrow1\).

Let service requirements \(S\) have cumulative distribution function
\(G\), with

\[
0<m=\mathbb ES<\infty,\qquad
G(x)\sim c x^\beta\quad(x\downarrow0),\qquad c,\beta>0.
\tag{G1}
\]

For this model and service law, write

\[
\varepsilon=1-\rho,\qquad
\alpha=\frac{\beta}{\beta+1},\qquad
\kappa=\left(\frac{(\beta+1)m}{c}\right)^{1/(\beta+1)}.
\]

\textbf{Theorem 1 (general-service heavy-traffic law).} If \(X,E\) are
independent unit-mean exponential variables, then

\[
\varepsilon^\alpha A_\rho\Rightarrow
V=\kappa X^\alpha E^{1-\alpha}.
\tag{G2}
\]

For every real \(p>0\),

\[
\lim_{\rho\uparrow1}\mathbb E[(\varepsilon^\alpha A_\rho)^p]
=\kappa^p\Gamma(1+\alpha p)\Gamma(1+(1-\alpha)p).
\tag{G3}
\]

In particular,

\[
\mathbb EA_\rho\sim
\frac{\kappa\Gamma(1+\alpha)\Gamma(2-\alpha)}{(1-\rho)^\alpha}.
\tag{G4}
\]

No finite second moment of \(S\) is assumed. The hypothesis constrains
only the lower end of the service distribution and its mean. A positive
lower bound on service requirements is outside the theorem.

The limiting distribution itself is a familiar product of powers of
exponential variables. The research claim is its occurrence as the
stationary age limit, with an exponent determined by the lower endpoint
of \(G\).

\section{3. Proof of the general-service
theorem}\label{proof-of-the-general-service-theorem}

\subsection{3.1 Normalisation and classical stationary
input}\label{normalisation-and-classical-stationary-input}

Replace clock time by \(t/m\) and service requirements by \(S/m\). The
resulting mean service requirement is one, arrival rate is \(\rho\), and
service CDF is

\[
F(x)=G(mx)\sim a x^\beta,\qquad a=c m^\beta.
\tag{G5}
\]

We prove the result in these units and restore \(m\) at the end.

The classical insensitivity theorem for the M/G/1 processor-sharing
queue gives

\[
\mathbb P(N=n)=(1-\rho)\rho^n\quad(n\geq0).
\tag{G6}
\]

Conditional on \(N=n\), the residual service requirements have the
product equilibrium-residual distribution. These facts follow, for
example, from Zachary (2007), Theorem 1: take his arrival function
\(\alpha(n)=\rho\), aggregate service function \(\beta(n)=1\) for
\(n>0\), and unit-mean service law \(F\). His balance equations reduce
to \(\pi(n+1)=\rho\pi(n)\). The input law here has no atom at zero by
(G1), and \(\sum_n\pi(n)\alpha(n)=\rho<\infty\).

The only parts of this stationary input needed for the proof are (G6),
finiteness of \(\mathbb EN\), and stationarity. If \(D_o(t)\) is the
total number of departures from an ordinary stationary queue during
\([0,t]\), conservation of customers gives

\[
\mathbb ED_o(t)=\rho t.
\tag{G7}
\]

Indeed, \(D_o(t)=N(0)+B(t)-N(t)\), the expectations of the two
stationary populations cancel, and the mean arrival count is \(\rho t\).
No second service moment or bounded departure intensity is needed.

\subsection{3.2 An exact auxiliary
process}\label{an-exact-auxiliary-process}

Initialise the queue at stationarity at zero. Call its \(N\) initial
packets old, and call every packet arriving after zero fresh. For fixed
\(t>0\), the event \(A_\rho(0)>t\) says that no packet generated after
\(-t\) has completed by zero. Shifting this event forward by \(t\) and
using stationarity gives

\[
\mathbb P(A_\rho>t)
=\mathbb P(\text{no fresh packet completes during }[0,t]).
\tag{G8}
\]

Construct an auxiliary queue from the same initial old residuals and the
same Poisson arrival epochs. Every fresh packet remains forever and
continues sharing the processor, irrespective of its finite service
requirement. Old packets retain their original residual requirements and
depart when those are exhausted. Let \(B(t)\) be the fresh-arrival
count, \(M(t)\) the auxiliary old population, and

\[
Q^*(t)=M(t)+B(t).
\]

The auxiliary evolution depends on the old residuals and the fresh
arrival epochs. It is independent of every fresh packet's service
requirement. A packet arriving at \(a_j\in(0,t]\) receives auxiliary
service

\[
V_j(t)=\int_{a_j}^t\frac{du}{Q^*(u)}.
\tag{G9}
\]

The denominator is positive on this integration interval because the
fresh packet itself remains present. The ordinary queue and the
auxiliary queue coincide until the first fresh completion. Hence no
fresh completion occurs by \(t\) exactly when every fresh requirement
exceeds its auxiliary service. Conditional independence of the fresh
requirements therefore yields the exact identity

\[
\mathbb P(A_\rho>t)
=\mathbb E\prod_{j:a_j\leq t}\bigl(1-F(V_j(t))\bigr).
\tag{G10}
\]

The empty product is one. This formula is valid for general service
CDFs, including distributions with atoms away from zero.

\subsection{3.3 Why old departures stay
negligible}\label{why-old-departures-stay-negligible}

Couple the ordinary queue used in (G7) to the auxiliary queue using the
same arrivals and finite service requirements, except that auxiliary
fresh requirements are infinite. On every finite interval there are
finitely many arrivals and initially finitely many packets, so the
following comparison can be established by induction through arrival and
departure events.

At every time, each packet still present in the ordinary queue is also
present in the auxiliary queue. For any common old packet, its ordinary
residual work is at most its auxiliary residual work. The assertion
holds initially. Between events, the ordinary queue has at most as many
packets, so it gives each common packet at least as much service. An
ordinary departure preserves the inclusion. If an auxiliary old packet
completes, its ordinary counterpart has already completed or completes
at the same time, so inclusion is again preserved. Fresh packets never
leave the auxiliary system. A common arrival preserves the assertion.

Consequently every auxiliary old departure occurs no earlier than the
corresponding ordinary departure, and

\[
0\leq N-M(t)\leq D_o(t).
\tag{G11}
\]

Take \(t=s\varepsilon^{-\alpha}\). Since \(\alpha<1\), (G7) and Markov's
inequality imply

\[
\varepsilon D_o(s\varepsilon^{-\alpha})\to0
\quad\text{in probability},
\qquad
\varepsilon B(s\varepsilon^{-\alpha})\to0
\quad\text{in probability}.
\tag{G12}
\]

Together with (G6), this proves that the scaled auxiliary population is
frozen on each compact rescaled time interval:

\[
\varepsilon N\Rightarrow X\sim\operatorname{Exp}(1),\qquad
\sup_{0\leq u\leq s}
|\varepsilon Q^*(u\varepsilon^{-\alpha})-\varepsilon N|
\to0\quad\text{in probability}.
\tag{G13}
\]

This is the step for which exponential service would have been
unnecessarily restrictive. The departure-count identity (G7) supplies
the needed bound for every finite-mean service distribution.

\subsection{3.4 Convergence of the exact
product}\label{convergence-of-the-exact-product}

Let \(u_j=\varepsilon^\alpha a_j\), and put
\(t=s\varepsilon^{-\alpha}\). From (G9),

\[
V_j(t)=\varepsilon^{1-\alpha}
\int_{u_j}^s\frac{du}{\varepsilon Q^*(u\varepsilon^{-\alpha})}.
\tag{G14}
\]

After restricting to \(\varepsilon N\geq\delta>0\), (G13) shows,
uniformly over all \(j\) with \(u_j\leq s\), that the last integral is

\[
\frac{s-u_j}{\varepsilon N}+o_{\mathbb P}(1).
\tag{G15}
\]

It is useful to make the summation control explicit. On a
high-probability set on which \(\varepsilon Q^*\) is between
\(\varepsilon N-\eta\) and \(\varepsilon N+\eta\), all the service
integrals in (G14) lie between

\[
\frac{\varepsilon^{1-\alpha}(s-u_j)}{\varepsilon N+\eta}
\quad\text{and}\quad
\frac{\varepsilon^{1-\alpha}(s-u_j)}{\varepsilon N-\eta}.
\]

This multiplicative sandwich avoids accumulating a separate
approximation error for every arrival. The Poisson law of large numbers
gives

\[
\varepsilon^\alpha\sum_{j:u_j\leq s}(s-u_j)^\beta
\longrightarrow\int_0^s(s-u)^\beta\,du
=\frac{s^{\beta+1}}{\beta+1}
\quad\text{in probability}.
\tag{G16}
\]

Its expectation tends to the displayed integral because the rescaled
arrival intensity is \(\rho\varepsilon^{-\alpha}\); its variance is
\(O(\varepsilon^\alpha)\). As \(\beta(1-\alpha)=\alpha\), (G5),
(G14)-(G16), and then \(\eta\downarrow0\), yield

\[
\sum_{j:a_j\leq t}F(V_j(t))
\Rightarrow \frac{a s^{\beta+1}}{(\beta+1)X^\beta}.
\tag{G17}
\]

On the same high-probability sets, \(\max_j V_j(t)\to0\), so
\(\max_j F(V_j(t))\to0\). Thus

\[
\left|\sum_j\log(1-F(V_j(t)))+\sum_jF(V_j(t))\right|
\leq\frac{\max_jF(V_j(t))}{1-\max_jF(V_j(t))}\sum_jF(V_j(t))
\to0
\]

in probability. Finally remove the restriction
\(\varepsilon N\geq\delta\) by letting \(\delta\downarrow0\), since its
limiting excluded probability is \(1-e^{-\delta}\).

The product in (G10) is bounded between zero and one. Its expectation
therefore converges, and

\[
\mathbb P(\varepsilon^\alpha A_\rho>s)
\longrightarrow
\mathbb E\exp\left\{-\frac{a s^{\beta+1}}{(\beta+1)X^\beta}\right\}.
\tag{G18}
\]

This is the survival function of

\[
\left(\frac{\beta+1}{a}\right)^{1/(\beta+1)}X^\alpha E^{1-\alpha}
\]

in mean-one service units.

\subsection{3.5 Uniform exponential tails and moment
convergence}\label{uniform-exponential-tails-and-moment-convergence}

Choose \(x_0>0\) and \(a_0>0\) such that

\[
F(x)\geq a_0x^\beta\qquad(0\leq x\leq x_0).
\tag{G19}
\]

For sufficiently small \(\varepsilon\), set \(t=s\varepsilon^{-\alpha}\)
and consider the event

\[
\mathcal H=\{N\leq s/\varepsilon,\ B(t)\leq2\rho t,
\ B(t/2)\geq\rho t/4\}.
\]

The geometric and Poisson Chernoff bounds imply, uniformly over
\(\rho\geq1/2\),

\[
\mathbb P(\mathcal H^c)
\leq e^{-s}+e^{-\rho t/3}+e^{-\rho t/16}
\leq e^{-s}+e^{-s/6}+e^{-s/32}.
\tag{G20}
\]

Since \(Q^*(u)\leq N+B(t)\) for \(u\leq t\), every arrival before
\(t/2\) receives auxiliary service at least

\[
\frac{t}{2(N+B(t))}
\geq\frac{\varepsilon^{1-\alpha}}{2(1+2\rho\varepsilon^{1-\alpha})}
\geq\frac16\varepsilon^{1-\alpha}=:r_\varepsilon
\quad\text{on }\mathcal H.
\tag{G21}
\]

Choose \(\varepsilon\) small enough that \(r_\varepsilon\leq x_0\). On
\(\mathcal H\), (G10) and (G19)-(G21) bound the conditional
no-completion probability by

\[
\exp\{-B(t/2)F(r_\varepsilon)\}
\leq\exp\left\{-\frac{\rho s\varepsilon^{-\alpha}}4
\frac{a_0\varepsilon^{\beta(1-\alpha)}}{6^\beta}\right\}
\leq\exp\left\{-\frac{a_0s}{8\,6^\beta}\right\}.
\tag{G22}
\]

Thus there are constants \(C,d>0\), depending only on \(F\), such that

\[
\mathbb P(\varepsilon^\alpha A_\rho>s)\leq C e^{-ds}
\tag{G23}
\]

for all \(s>0\) and all sufficiently small \(\varepsilon\). Dominated
convergence in the tail-integral formula now proves all positive moment
limits. Restoring the original time units replaces the mean-one
prefactor by

\[
m\left(\frac{\beta+1}{c m^\beta}\right)^{1/(\beta+1)}
=\left(\frac{(\beta+1)m}{c}\right)^{1/(\beta+1)}=\kappa.
\]

This proves (G2)-(G4). \(\square\)

\subsection{3.6 Regularly varying service CDFs: a scale beyond pure
powers}\label{regularly-varying-service-cdfs-a-scale-beyond-pure-powers}

A pure-power approximation suppresses factors such as \(\log(1/x)\).
Such factors can change the leading age scale even though the power
exponent remains the same. The following extension gives a normalization
directly from the service CDF, without differentiability or a finite
second moment.

\textbf{Theorem 1a (regularly varying lower endpoint).} Keep the queue
model of Section 2 and the assumption \(0<m=\mathbb ES<\infty\). Suppose
the mean-one service CDF \(F(x)=G(mx)\) satisfies \[
\lim_{x\downarrow0}\frac{F(tx)}{F(x)}=t^\beta
\quad\text{for every }t>0,\qquad \beta>0. \tag{R1}
\] For sufficiently small \(\varepsilon>0\), define \(q_\varepsilon>0\)
uniquely by \[
\int_0^{q_\varepsilon}F(u)\,du=\frac{\varepsilon}{\beta+1},
\qquad b_\varepsilon=\frac{q_\varepsilon}{\varepsilon}. \tag{R2}
\] Then \(q_\varepsilon\to0\), \(b_\varepsilon\to\infty\), and \[
\frac{A_\rho}{m b_\varepsilon}\Rightarrow
Z_\beta=(\beta+1)^{1/(\beta+1)}
X^{\beta/(\beta+1)}E^{1/(\beta+1)}. \tag{R3}
\] For every \(p>0\), the \(p\)-th moment converges to \[
(\beta+1)^{p/(\beta+1)}
\Gamma\!\left(1+\frac{\beta p}{\beta+1}\right)
\Gamma\!\left(1+\frac{p}{\beta+1}\right). \tag{R4}
\]

\textbf{Proof.} Work in mean-one units. Assumption (R1) implies
\(F(0)=0\) and \(F(x)>0\) for all sufficiently small \(x>0\). Thus the
integral in (R2) is continuous and strictly increasing near zero, even
when \(F\) has jumps. The normalization is consequently well defined.

The elementary integral consequence of (R1) needed here follows by
changing variables: \[
\frac{\int_0^qF(u)\,du}{qF(q)}
=\int_0^1\frac{F(qv)}{F(q)}\,dv
\longrightarrow\int_0^1v^\beta\,dv=\frac1{\beta+1}. \tag{R5}
\] Dominated convergence applies because the ratio is between zero and
one. Hence \[
b_\varepsilon F(q_\varepsilon)\longrightarrow1,
\qquad q_\varepsilon=\varepsilon b_\varepsilon\longrightarrow0. \tag{R6}
\] Since \(F(q_\varepsilon)\to0\), this also proves
\(b_\varepsilon\to\infty\).

Use the same auxiliary queue and exact product (G10), now on the time
interval \(t=s b_\varepsilon\). Conservation and the monotone coupling
still give \[
\varepsilon B(s b_\varepsilon)\to0,\qquad
\varepsilon D_o(s b_\varepsilon)\to0
\] in probability, since their expectations are
\(\rho s q_\varepsilon\). Thus the scaled initial population is again
asymptotically frozen and converges to \(X\).

Fix \(0<\delta<K<\infty\), and restrict first to
\(\delta\leq\varepsilon N\leq K\). For a fixed sufficiently small
\(\eta>0\), the uniform population comparison bounds the service
exposure of an arrival at rescaled time \(u=a_j/b_\varepsilon\) between
\[
q_\varepsilon\frac{s-u}{\varepsilon N+\eta}
\quad\text{and}\quad
q_\varepsilon\frac{s-u}{\varepsilon N-\eta}.
\] Convergence in (R1) is uniform for multipliers in every compact
interval \([0,L]\). To see this without a supplementary regularity
assumption, use monotonicity of \(F(qv)/F(q)\), pointwise convergence
and a finite mesh for the continuous limit \(v^\beta\); the interval
next to zero is bounded by its right endpoint. The weighted Poisson law
of large numbers, (R6), and this uniform convergence therefore yield \[
\sum_{a_j\leq s b_\varepsilon}F(V_j(s b_\varepsilon))
\Rightarrow\frac{s^{\beta+1}}{(\beta+1)X^\beta}. \tag{R7}
\] More explicitly, after division by \(b_\varepsilon\) the arrival sum
of any bounded continuous weight converges to its integral; its variance
is \(O(b_\varepsilon^{-1})\). The multiplicative sandwich above,
followed by \(\eta\downarrow0\), proves (R7). The largest summand tends
to zero, so the logarithmic remainder estimate in Section 3.4 applies.
Remove the restriction by \(\delta\downarrow0\) and \(K\uparrow\infty\).
The bounded product then proves (R3).

For moments, use \[
\mathcal H=\{N\leq s/\varepsilon,\ B(t)\leq2\rho t,\
B(t/2)\geq\rho t/4\},\qquad t=s b_\varepsilon .
\] For small \(\varepsilon\), \(b_\varepsilon\geq1\), so the complement
has the same bound as (G20). Every arrival before \(t/2\) receives at
least \[
\frac{q_\varepsilon}{2(1+2\rho q_\varepsilon)}
\geq q_\varepsilon/6
\] on \(\mathcal H\). Assumptions (R1) and (R6) imply \[
b_\varepsilon F(q_\varepsilon/6)\longrightarrow6^{-\beta}>0.
\] The conditional survival probability is therefore at most
\(e^{-d_0s}\) for a constant \(d_0>0\) independent of sufficiently small
\(\varepsilon\). Adding the exceptional-event probabilities gives
\(\mathbb P(A_\rho/(m b_\varepsilon)>s)\leq C e^{-ds}\). Tail
integration proves (R4). \(\square\)

For \(F(x)\sim ax^\beta\), (R2) gives
\(b_\varepsilon\sim a^{-1/(\beta+1)}\varepsilon^{-\beta/(\beta+1)}\);
Theorem 1 is recovered exactly. More generally, if
\(F(x)\sim ax^\beta(\log(1/x))^\gamma\), with any real \(\gamma\), then
\[
b_\varepsilon\sim
\left(\frac{(\beta+1)^\gamma}{a}\right)^{1/(\beta+1)}
\varepsilon^{-\beta/(\beta+1)}
\bigl(\log(1/\varepsilon)\bigr)^{-\gamma/(\beta+1)}. \tag{R8}
\] Indeed, substitution into \(qF(q)\sim\varepsilon\) gives the
displayed constant because
\(\log(1/q)\sim\log(1/\varepsilon)/(\beta+1)\). This example is a
genuine enlargement of the pure-power statement: omitting the
logarithmic factor gives the wrong leading scale when \(\gamma\ne0\).

\section{4. Service-distribution
consequences}\label{service-distribution-consequences}

For Erlang service with \(k\) stages of rate \(k/m\),

\[
\beta=k,\qquad c=\frac{(k/m)^k}{k!},\qquad
\alpha=\frac{k}{k+1},\qquad
\kappa=m\left(\frac{(k+1)k!}{k^k}\right)^{1/(k+1)}.
\]

Thus the age exponents for one-, two- and three-stage Erlang services
are \(1/2,2/3,3/4\). All three queues have exactly the same stationary
population distribution at the same load and mean service time. Their
information-age divergence rates differ because their probabilities of
very short service requirements differ.

For a Weibull service distribution with shape \(\beta\), scale
\(\theta\), and mean \(m=\theta\Gamma(1+1/\beta)\), the lower-end
coefficient is \(c=\theta^{-\beta}\). The same exponent
\(\beta/(\beta+1)\) follows. The theorem allows any positive \(\beta\),
including \(\beta<1\).

For a finite mixture of exponential services with rates \(\mu_j\) and
probabilities \(p_j\),

\[
m=\sum_jp_j/\mu_j,\qquad c=\sum_jp_j\mu_j,\qquad\beta=1.
\]

At fixed mean \(m\), its leading mean-age coefficient is
\(\pi\sqrt{m/c}/(2\sqrt2)\). Since \(mc\geq1\), that coefficient is no
larger than the exponential-service coefficient with the same mean. This
is an asymptotic comparison of these service distributions; it is not a
general monotonicity theorem in the coefficient of variation.

\section{5. Exponential service, global bounds and the
conjecture}\label{exponential-service-global-bounds-and-the-conjecture}

\subsection{5.1 The complete exponential-service
limit}\label{the-complete-exponential-service-limit}

\textbf{Corollary 2.} For \(G(x)=1-e^{-\mu x}\), let
\(\varepsilon=1-\rho\). Then

\[
\mu\sqrt\varepsilon A_\rho\Rightarrow W:=\sqrt{2E_1E_2},
\qquad
\lim_{\rho\uparrow1}\mathbb E[(\mu\sqrt\varepsilon A_\rho)^p]
=2^{p/2}\Gamma(1+p/2)^2\quad(p>0),
\tag{M1}
\]

where \(E_1,E_2\) are independent unit-mean exponentials. Consequently,

\[
\mathbb EA_\rho\sim\frac{\pi}{2\sqrt2\mu\sqrt{1-\rho}},
\qquad
\operatorname{Var}(A_\rho)\sim
\frac{2-\pi^2/8}{\mu^2(1-\rho)}.
\tag{M2}
\]

The limiting survival function can be written

\[
\mathbb P(W>s)=\int_0^\infty e^{-x-s^2/(2x)}\,dx
=\sqrt2s K_1(\sqrt2s),\qquad s>0,
\tag{M3}
\]

with value one at zero. Here \(K_1\) is the modified Bessel function of
the second kind; the integral is its standard defining integral
representation after rescaling. The distribution is classical. All the
queueing statements follow by substituting \(m=1/\mu\), \(c=\mu\), and
\(\beta=1\) in Theorem 1.

\subsection{5.2 Bounds at every load}\label{bounds-at-every-load}

\textbf{Proposition 3.} For every \(0<\rho=\lambda/\mu<1\),

\[
\frac1\lambda+\frac1\mu
\leq\mathbb EA_\rho
\leq\frac1\mu\left(1+\frac1\rho+\frac{\rho^2}{1-\rho}\right).
\tag{M4}
\]

For the lower bound, use the forward representation. Let \(S\) be the
first arrival after zero. It has mean \(1/\lambda\). Independently,
potential service events form a Poisson process of rate \(\mu\), thinned
to remove idle events. The first fresh completion cannot precede the
first potential service event after \(S\). That residual event time has
mean \(1/\mu\). In fact this construction proves the stochastic lower
bound

\[
\operatorname{Exp}(\lambda)+\operatorname{Exp}(\mu)
\ \leq_{\mathrm{st}}\ A_\rho,
\tag{M5}
\]

with independent exponentials on the left.

For the upper bound, couple processor sharing and
first-come-first-served (FCFS) through common arrival times and common
rate-\(\mu\) potential service times. Under either discipline every
nonempty queue completes a packet at rate \(\mu\). Their queue lengths
and actual departure epochs therefore agree. Label arrivals after any
common empty epoch in increasing order. After \(k\) subsequent
completions, FCFS has delivered labels \(1,\ldots,k\). Processor sharing
has delivered \(k\) distinct labels, so its largest delivered label is
at least \(k\). Hence the freshest timestamp delivered by processor
sharing is at least as large at every time. At an empty epoch both
queues have delivered all earlier arrivals, so their freshest timestamps
agree. The positive recurrent birth-death chain has a past empty epoch
almost surely, which makes the same coupling valid in stationarity.
Therefore

\[
A_\rho\leq_{\mathrm{st}} A_\rho^{\mathrm{FCFS}}.
\tag{M6}
\]

For completeness, the classical FCFS mean can be obtained directly. Let
\(Y\sim\operatorname{Exp}(\lambda)\) be an interarrival time,
\(T_-\sim\operatorname{Exp}(\mu-\lambda)\) the preceding packet's
stationary system time, and \(S\sim\operatorname{Exp}(\mu)\) the current
packet's service. Then \(Y\) is independent of \(T_-\), and the current
system time is \(T=(T_--Y)^++S\). The exponential law of \(T_-\) follows
by mixing sums of geometric-many independent rate-\(\mu\) exponentials.
The usual area identity between consecutive delivery epochs yields

\[
\mathbb EA_\rho^{\mathrm{FCFS}}
=\frac1\lambda+\lambda\mathbb E[YT].
\]

Now \(\lambda\mathbb E[YS]=1/\mu\), while

\[
\mathbb E[Y(T_--Y)^+]
=\int_0^\infty y\lambda e^{-\lambda y}
\frac{e^{-(\mu-\lambda)y}}{\mu-\lambda}\,dy
=\frac{\lambda}{\mu^2(\mu-\lambda)}.
\]

Combining this computation with the stochastic comparison proves (M4).

Define exactly the function used in the source,

\[
C(\rho)=\mu\mathbb EA_\rho-1-\rho^{-1}.
\]

Then (M4) proves \(0\leq C(\rho)\leq\rho^2/(1-\rho)\) for every
\(\rho\). In particular \(C(\rho)\to0\) at zero. Corollary 2 gives

\[
\sqrt{1-\rho}\,C(\rho)\longrightarrow\frac\pi{2\sqrt2},
\tag{M7}
\]

so \(C(\rho)\to\infty\) at one. These statements establish all parts of
Conjecture 1 preceding equation (20).

Equation (20) in the inspected arXiv manuscript asserts, for
sufficiently large \(\rho\),

\[
\frac{(\rho-1/2)^3}{1-\rho}\leq C(\rho)
\leq\frac{3\rho}{4\sqrt{1-\rho}}.
\tag{M8}
\]

These two bounds are mutually incompatible even before the queue is
analysed: their ratio tends to infinity. The new limit also determines
which individual asymptotic claims fail. The left member grows too
quickly, while the right member has limiting coefficient
\(3/4<\pi/(2\sqrt2)\). Thus both sides of (M8) fail sufficiently close
to saturation. This attribution is to the inspected arXiv version.
Direct inspection of the journal version of record and any formal
correction remains a bibliographic limitation. The substantive theorem
does not depend on the printed inconsistency.

\section{6. Exact representations and reproducible
checks}\label{exact-representations-and-reproducible-checks}

\subsection{6.1 A two-coordinate generator for exponential
service}\label{a-two-coordinate-generator-for-exponential-service}

The stationary-age survival event is a first-fresh-completion event by
(G8). With exponential service, conditional on observable queue history
and attained services, unobserved residual requirements remain
independent exponentials. Before the first fresh completion, let \(m,n\)
be the old and fresh populations. The exact transient transitions are

\[
(m,n)\longrightarrow(m,n+1)\quad\text{at rate }\lambda,
\qquad
(m,n)\longrightarrow(m-1,n)\quad\text{at rate }\mu\frac{m}{m+n},
\]

with killing at rate \(\mu n/(m+n)\); both service rates are zero at the
empty state. Initialise \(m\) geometrically with
\(\mathbb P(m=k)=(1-\rho)\rho^k\) and \(n=0\). Simulation until killing
gives an exact sample from the stationary age distribution. It needs
neither a vector of timestamps nor a burn-in period.

Equivalently, suppress killing, give the old population deaths at rate
\(\mu M/(M+B)\), and let fresh arrivals remain forever. Then

\[
\mathbb P(A_\rho>t)=\mathbb E\exp\left\{-\int_0^t\mu\frac{B(u)}{M(u)+B(u)}\,du\right\}.
\tag{V1}
\]

This is also obtained directly from (G10) by using an exponential
survival function and summing all fresh service exposures. Transition
rates are bounded by \(\lambda+\mu\), so the process is nonexplosive.

\subsection{6.2 Finite-state exact
checks}\label{finite-state-exact-checks}

The accompanying \texttt{verify.py} solves the finite-capacity killed
chains by Gaussian elimination over rational numbers. At
\(\lambda=1/2\), \(\mu=1\), capacities one, two and three give the exact
mean ages

\[
\frac{10}{3},\qquad\frac{45}{14},\qquad\frac{1292}{405}.
\]

For capacity two the result exactly matches Proposition 1 of Gandarias
et al.~(2023). The finite-state check tests the first-completion
reduction and the handling of obsolete packets. The main theorem uses an
infinite buffer throughout and does not approximate it by a fixed
capacity.

\subsection{6.3 Independent stationary
samples}\label{independent-stationary-samples}

For exponential mean-one service, \texttt{code/verify.py} uses 200,000
independent samples at each listed load. The supplied research run's
full output, including tail frequencies and seeds, is retained in
\texttt{evidence/historical/verification.json}.

\Needspace{19\baselineskip}

{\def\LTcaptype{none} % do not increment counter
\begin{longtable}[]{@{}
  >{\raggedright\arraybackslash}p{(\linewidth - 4\tabcolsep) * \real{0.2727}}
  >{\raggedleft\arraybackslash}p{(\linewidth - 4\tabcolsep) * \real{0.3636}}
  >{\raggedleft\arraybackslash}p{(\linewidth - 4\tabcolsep) * \real{0.3636}}@{}}
\toprule\noalign{}
\begin{minipage}[b]{\linewidth}\raggedright
Load \(\rho\)
\end{minipage} & \begin{minipage}[b]{\linewidth}\raggedleft
Mean of \(\sqrt{1-\rho}A_\rho\)
\end{minipage} & \begin{minipage}[b]{\linewidth}\raggedleft
Monte Carlo standard error
\end{minipage} \\
\midrule\noalign{}
\endhead
\bottomrule\noalign{}
\endlastfoot
0.9000 & 1.285566 & 0.002112 \\
0.9900 & 1.145599 & 0.001999 \\
0.9990 & 1.116731 & 0.001965 \\
0.9999 & 1.110625 & 0.001962 \\
Limit & \(\pi/(2\sqrt2)=1.110721\) & --- \\
\end{longtable}
}

At \(\rho=0.9999\), the scaled second moment is 2.003068, compared with
the exact limit 2.

\texttt{verify\_erlang.py} checks non-exponential examples using an
exact phase-count chain. Conditional on the geometric stationary
population, the residual service of an Erlang-\(k\) packet is an equal
mixture of Erlang distributions with one through \(k\) phases. This
follows from the equilibrium-residual law in Section 3.1. Each phase has
rate \(k\) in mean-one units. A new packet enters with \(k\) phases; a
fresh packet's final phase completion kills the chain. Thus these too
are direct stationary-age samples, without burn-in.

\Needspace{10\baselineskip}

{\def\LTcaptype{none} % do not increment counter
\begin{longtable}[]{@{}
  >{\raggedright\arraybackslash}p{(\linewidth - 12\tabcolsep) * \real{0.1111}}
  >{\raggedleft\arraybackslash}p{(\linewidth - 12\tabcolsep) * \real{0.1481}}
  >{\raggedleft\arraybackslash}p{(\linewidth - 12\tabcolsep) * \real{0.1481}}
  >{\raggedleft\arraybackslash}p{(\linewidth - 12\tabcolsep) * \real{0.1481}}
  >{\raggedleft\arraybackslash}p{(\linewidth - 12\tabcolsep) * \real{0.1481}}
  >{\raggedleft\arraybackslash}p{(\linewidth - 12\tabcolsep) * \real{0.1481}}
  >{\raggedleft\arraybackslash}p{(\linewidth - 12\tabcolsep) * \real{0.1481}}@{}}
\toprule\noalign{}
\begin{minipage}[b]{\linewidth}\raggedright
Erlang shape
\end{minipage} & \begin{minipage}[b]{\linewidth}\raggedleft
Age exponent
\end{minipage} & \begin{minipage}[b]{\linewidth}\raggedleft
Samples per load
\end{minipage} & \begin{minipage}[b]{\linewidth}\raggedleft
Scaled mean at 0.99
\end{minipage} & \begin{minipage}[b]{\linewidth}\raggedleft
At 0.999
\end{minipage} & \begin{minipage}[b]{\linewidth}\raggedleft
At 0.9999
\end{minipage} & \begin{minipage}[b]{\linewidth}\raggedleft
Limit
\end{minipage} \\
\midrule\noalign{}
\endhead
\bottomrule\noalign{}
\endlastfoot
2 & \(2/3\) & 100,000 & 1.006968 & 0.961269 & 0.942002 & 0.922792 \\
3 & \(3/4\) & 50,000 & 0.934616 & 0.869475 & 0.840009 & 0.808869 \\
\end{longtable}
}

At the last load the standard errors are respectively 0.002392 and
0.003112. These values show appreciable finite-load bias beyond
simulation uncertainty; the theorem asserts a limit and supplies no
convergence rate. The bias decreases in the displayed range. Complete
supplied results, including second moments, are in
\texttt{evidence/historical/erlang\_verification.json}. Computation
corroborates the reductions and asymptotic predictions; it is separate
from the proof.

The publication check \texttt{code/replay.py} independently reconstructs
the rational finite-state linear systems, checks the pure-power scale
algebra, tests a logarithmic scale example and runs 2,000 stationary
exponential samples with a new seed. It does not rerun the much larger
historical tables. Its recorded environment and results are supplied
separately. The historical runs did not record the exact original NumPy
version, so bitwise reproduction of their random streams is not
asserted.

\section{7. Interpretation, prior-art boundary and
limitations}\label{interpretation-prior-art-boundary-and-limitations}

The initial population is of order \(\varepsilon^{-1}\), with an
asymptotically exponential scaled size. During an interval of length
\(\varepsilon^{-\alpha}\), each early fresh packet receives work of
order \(\varepsilon^{1-\alpha}\), and the number of fresh packets is of
order \(\varepsilon^{-\alpha}\). The small-service assumption gives an
individual completion probability of order
\(\varepsilon^{\beta(1-\alpha)}\). Balancing this against the number of
fresh packets gives \(\beta(1-\alpha)=\alpha\), or
\(\alpha=\beta/(\beta+1)\). The integral over arrival times and the
random initial population supply the exact distribution and constant.
The proof controls both the changing backlog and the large-age tail,
which are needed to justify this scaling argument and passage to all
moments.

At equal load and equal mean service, the stationary population is
identical for exponential, Erlang and Weibull service. Their age
exponents can differ. This provides a specific limitation of using
ordinary queue occupancy as a proxy for freshness. For exponential
service, the FCFS mean age grows at order \((1-\rho)^{-1}\), whereas the
processor-sharing mean age grows at order \((1-\rho)^{-1/2}\). Both
still diverge as saturation is approached.

The literature search inspected the originating 2023 conjecture, the
2021 scheduling study, primary insensitivity theorems, subsequent age
papers on shared computing systems, and formula-level aliases for the
limiting law and first-completion construction. Chiariotti (2024) uses
periodic requests and cancellation/replacement of old requests for each
client; its model differs from the present lossless infinite-buffer
queue. Inoue and Mandjes (2025) analyse a shared server under FCFS,
identified by their workload-to-sojourn formula, rather than egalitarian
processor sharing. Redder (2026) studies parallel renewal processes
rather than this queueing model. No equivalent general-service or
exponential heavy-traffic age theorem was located in the bounded corpus
recorded in \texttt{novelty\_report.md}. This is not a proof of
historical priority or a claim that no later resolution exists.

The assumptions have concrete limits. The arrival process is Poisson,
scheduling is equal processor sharing, obsolete jobs remain until
service completion, and the monitor retains the maximum delivered
generation time. A positive minimum service requirement, an atom at
zero, time-varying service laws, deterministic sampling, finite buffers,
and cancellation policies are outside both theorems. Theorem 1a includes
general regularly varying CDFs of positive index, but does not include
index zero or rapidly varying lower endpoints. Neither theorem gives a
closed-form finite-load mean or a numerical error bound for a
heavy-traffic approximation. The global stochastic comparisons in
Section 5 use exponential service and are not asserted for arbitrary
\(G\).

The original pure-power proof underwent a separate AI adversarial review
within the research workflow. Such review is not external human peer
review. The theorem's provenance and significance remain subject to the
documented literature-search boundary and to normal scholarly
assessment.

\section{8. Research declarations}\label{research-declarations}

\textbf{Data and code availability.} The accompanying repository
contains this complete manuscript, the finite-state and simulation code,
historical numerical outputs, a bounded fresh replay, a
claim-to-evidence map and a rating-free revision response. The source
repository is
\url{https://github.com/ipitchford/processor-sharing-freshness} and the
versioned archive is \url{https://doi.org/10.5281/zenodo.23501620}. The
mathematical proofs are self-contained apart from the explicitly cited
classical stationary-queue theorem. No private review or inaccessible
development note is needed to follow the argument.

\textbf{Ethics.} This theoretical and computational study uses no human
participants, personal data or animal subjects.

\textbf{Contributions.} The release is attributed to Anonymous. The
public evidence records the mathematical development, software and
validation scopes. It does not assign unverified individual human CRediT
roles.

\textbf{Funding and competing interests.} No external research funding
or competing interest has been declared in the supplied materials.

\textbf{AI assistance.} AI agents assisted with mathematical
development, literature checking, code, manuscript preparation and
adversarial review. The accompanying computational results are
distinguished from analytic proof. AI review is not represented as
independent human peer review or formal proof verification.

\section{References}\label{references}

Bonald, T., \& Proutière, A. (2002). Insensitivity in processor-sharing
networks. \emph{Performance Evaluation, 49}, 193--209.
\url{https://doi.org/10.1016/S0166-5316(02)00110-4}

Chiariotti, F. (2024). Age of information analysis for a shared edge
computing server. \emph{IEEE Transactions on Communications, 72}(12),
7826--7841. \url{https://doi.org/10.1109/TCOMM.2024.3412761}

Gandarias, B., Doncel, J., \& Assaad, M. (2023). On the age of
information of processor sharing systems. \emph{Journal of
Communications and Networks, 25}(5), 643--656.
\url{https://doi.org/10.23919/JCN.2023.000042} . Inspected manuscript:
\url{https://arxiv.org/pdf/2309.02083} . The conjecture attribution and
fake-update convention refer specifically to Conjecture 1, equations
(19)--(20), and Appendix A, Table I, of that version.

Inoue, Y., \& Mandjes, M. (2025). Characterizing the age of information
with multiple coexisting data streams. \emph{IEEE Transactions on
Information Theory, 71}(6), 4732--4753.
\url{https://doi.org/10.1109/TIT.2025.3544266}

Liu, Z., Huang, L., Li, B., \& Ji, B. (2021). Anti-aging scheduling in
single-server queues: A systematic and comparative study. \emph{Journal
of Communications and Networks, 23}(2), 91--105.
\url{https://doi.org/10.23919/JCN.2021.000005} . Inspected manuscript:
\url{https://arxiv.org/pdf/2003.04271}

Redder, A. (2026). Age-of-information in distributed systems caused by
asynchronous computing modeled as parallel renewal processes.
\emph{Statistics and Computing, 36}, Article 136.
\url{https://doi.org/10.1007/s11222-026-10885-x}

Zachary, S. (2007). A note on insensitivity in stochastic networks.
\emph{Journal of Applied Probability, 44}(1), 238--248.
\url{https://doi.org/10.1239/jap/1175267175} . Theorem 1; preprint:
\url{https://arxiv.org/abs/math/0611526}

\end{document}
